% Chương 1
\chapter{Giới thiệu chung}
\label{ch:chuong1}

\section{Đặt vấn đề}

\textbf{Bối cảnh:} Da liễu là chuyên khoa có số lượng bệnh lý lớn, biểu hiện đa dạng và mang tính trực quan cao. Do rào cản tiếp cận bác sĩ chuyên khoa, người bệnh thường tự tìm kiếm thông tin trên internet trước khi thăm khám.

\textbf{Khó khăn thực tế:}
\begin{itemize}
  \item Người bệnh: Gặp rào cản lệch từ vựng (mô tả đời thường so với thuật ngữ y khoa), dễ nhầm lẫn giữa các bệnh có hình thái tương tự, tiếp cận nguồn thông tin kém kiểm chứng và đối mặt với rủi ro tự điều trị sai cách.
  \item Bác sĩ: Tốn thời gian khai thác lại bệnh sử chưa được chuẩn hóa từ đầu trong mỗi buổi khám.
\end{itemize}

\textbf{Nhu cầu:} Xây dựng công cụ hỗ trợ tiền chẩn đoán đáng tin cậy với các đặc tính kỹ thuật:
\begin{itemize}
  \item Tiếp nhận và xử lý đa phương thức (mô tả bằng văn bản và ảnh chụp vùng da).
  \item Cung cấp thông tin tham khảo có nguồn trích dẫn kiểm chứng.
  \item Chủ động khai thác thông tin thiếu hụt như quy trình hỏi bệnh sử lâm sàng.
  \item Tuân thủ giới hạn y khoa (không đưa ra chẩn đoán xác định, không kê đơn và cảnh báo thời điểm cần gặp bác sĩ).
\end{itemize}

\section{Phạm vi}

\section{Đề xuất giải pháp}

\section{Kết luận chương}

Chương này đã nêu bối cảnh của đề tài: da liễu là chuyên khoa có nhiều bệnh lý, biểu hiện đa dạng và mang tính trực quan cao, trong khi người bệnh gặp rào cản về từ vựng và dễ nhầm lẫn giữa các bệnh tương tự, còn bác sĩ tốn thời gian khai thác lại bệnh sử chưa được chuẩn hóa. Từ đó, chương xác định nhu cầu về một công cụ hỗ trợ tiền chẩn đoán đáng tin cậy với bốn đặc tính: xử lý đa phương thức (văn bản và ảnh), cung cấp thông tin tham khảo có nguồn trích dẫn, chủ động khai thác thông tin còn thiếu như quy trình hỏi bệnh sử, và tuân thủ giới hạn y khoa.

Các yêu cầu trên là cơ sở để đề tài kết hợp mạng CNN phân loại ảnh, truy xuất tăng cường sinh văn bản (RAG) và đồ thị tri thức, được trình bày chi tiết trong các chương tiếp theo.
